\chapter{Case studies: five engines, told in operations language}
\label{ch:cases}

\begin{quote}\itshape
This chapter in brief: five real engines from the test fleet, each illustrating one
situation a fleet manager actually faces, told with the monitoring plots in front of us and
both pipelines' answers on record. Every engine, number and verdict comes from
\texttt{scripts/make\_case\_studies.py} and the tuned F5 configurations --- nothing staged.
\end{quote}

\keyidea{Told in operations language, the same findings become actionable: prognosis buys a plannable shop-visit slot, the confusable wall is a purchase order to a sensor vendor, and a drifting thermocouple corrupts every downstream method --- sensor hygiene is monitoring infrastructure.}


\section*{Case A --- Engine 17: ordinary aging, and why wash scheduling works}

Engine 17 lives 18\,323 cycles with no acute fault and no sensor trouble: the baseline
customer of any monitoring system. Its EGT margin (\cref{fig:case-a}) shows textbook fleet
physics --- fast early erosion, sixteen wash recoveries, the long slow slide to retirement.
Operationally, this is the case where the \emph{traditional} toolkit earns its keep at
near-zero cost: margin trending plus the wash log explains everything visible, and nothing
here needs a neural network. A monitoring roadmap should leave this workflow alone.

\begin{figure}[htbp]
  \centering
  \includegraphics[width=0.78\textwidth]{case_a_fouling}
  \caption{Case A: EGT margin of engine 17. Gray verticals: compressor washes. Auto-generated by \texttt{scripts/make\_case\_studies.py}.}
  \label{fig:case-a}
\end{figure}

\section*{Case B --- Engine 89: the hot-section burner, where prognosis pays}

Engine 89 drew a severe hot-section multiplier (2.0$\times$ fleet nominal) and retires at
11\,470 cycles --- the short-life tail every planner fears. At the half-life review the
tuned margin extrapolation was telling maintenance control ``about 6\,300 cycles more than
reality''; the GRU's error at the same moment was 2\,234 cycles, tightening to 1\,396 by the
90\,\% point (\cref{fig:case-b}). Six thousand optimistic cycles is the difference between
an orderly shop-visit slot and a schedule scramble; this single engine is the H3 verdict
made concrete.

\begin{figure}[htbp]
  \centering
  \includegraphics[width=0.78\textwidth]{case_b_prognosis}
  \caption{Case B: EGT margin of engine 89 with the three review points (50/70/90\,\% of
  life) marked. Auto-generated by \texttt{scripts/make\_case\_studies.py}.}
  \label{fig:case-b}
\end{figure}

\section*{Case C --- Engine 0: the lying thermocouple}

Engine 0 carries a drifting EGT sensor that accumulates a $+9$\,K bias by end of life
(\cref{fig:case-c}). Every deviation-based consumer downstream --- trending, GPA, margin
estimation --- ingests that bias as if the gas path were 9\,K hotter than it is: a
phantom deterioration of roughly a third of a typical year's real margin loss. The classical
remedy --- an augmented-state Kalman that estimates the bias (\cref{sec:gpa}) --- is built
and tested in \cref{sec:f8-l7}: it \emph{tracks} the drift well but, on the cockpit rank-3
set, barely un-corrupts the diagnosis, so the operational lesson survives any modeling
choice: \textbf{cross-channel consistency checks and periodic probe calibration are
monitoring infrastructure, not optional hygiene} --- the smartest algorithm inherits its
sensors' honesty.

\begin{figure}[htbp]
  \centering
  \includegraphics[width=0.78\textwidth]{case_c_drift}
  \caption{Case C: measured-minus-true EGT of engine 0 (smoothed): the drift ramp the
  monitoring system cannot see from this channel alone. Auto-generated by \texttt{scripts/make\_case\_studies.py}.}
  \label{fig:case-c}
\end{figure}

\section*{Case D --- Engine 13: bird at cycle 5\,727, alarm at 5\,749}

A foreign-object event costs engine 13 about 1.4\,\% of compressor efficiency in one
flight. The cruise EGT deviation steps visibly (\cref{fig:case-d}) and the tuned traditional
CUSUM --- fed by Holt innovations, per the chapter-6 lesson --- alarms \textbf{22 cycles}
after the event: for a daily utilization of 4--6 cycles, that is a borescope request within
the week. Step events are the classical detector's home turf and this case shows the
industrial stack at its honest best.

\begin{figure}[htbp]
  \centering
  \includegraphics[width=0.78\textwidth]{case_d_fod}
  \caption{Case D: engine 13's cruise EGT deviation (smoothed), the FOD event (dashed) and
  the traditional alarm (dotted), 22 cycles later. Auto-generated by \texttt{scripts/make\_case\_studies.py}.}
  \label{fig:case-d}
\end{figure}

\section*{Case E --- Engine 0 again: the confusable fault, live}

Late in life (cycle 14\,061) engine 0 develops a genuine HPC-efficiency fault --- one pole
of the 1.3$^{\circ}$ confusable pair. Asked at onset$+500$, the tuned ICM rule answers
\emph{hpc.eta}: correct. The tuned AI classifier answers \emph{hpt.eta} --- precisely the
confusable partner (\cref{fig:case-e}). On the earlier, subtler LPC episode of the same
engine, both families miss in different directions. This is the H2 refutation embodied: at
these signature angles the cockpit channels genuinely cannot separate the pair, and which
family lands the point is close to coin-flip territory. The purchase order that fixes this
case is written to a sensor vendor, not a software one.

\begin{figure}[htbp]
  \centering
  \includegraphics[width=0.78\textwidth]{case_e_confusable}
  \caption{Case E: ground truth of the faulted parameter and its confusable partner on
  engine 0 (onset dashed). The AI attributed the step to the partner. Auto-generated by \texttt{scripts/make\_case\_studies.py}.}
  \label{fig:case-e}
\end{figure}
